\documentclass[11pt,a4paper]{article}

\usepackage{iftex}
\ifPDFTeX
  \usepackage[utf8]{inputenc}
  \usepackage[T1]{fontenc}
  \usepackage{lmodern}
\else
  \usepackage{fontspec}
\fi
\usepackage[brazilian]{babel}
\usepackage[margin=2.5cm]{geometry}
\usepackage{booktabs}
\usepackage{float}
\usepackage{graphicx}
\usepackage{amsmath}
\usepackage{enumitem}
\usepackage{xcolor}
\usepackage{tikz}
\usetikzlibrary{positioning,arrows.meta}
\usepackage[hidelinks]{hyperref}

\newcommand{\cfg}[1]{\texttt{\detokenize{#1}}}
\setlist{itemsep=2pt,topsep=4pt}

\title{Q1: comparação de documentos semelhantes no mesmo caso\\[4pt]
\large Cópia, versão modificada ou documento diferente}
\author{Guilherme --- Case técnico Enter}
\date{16 de outubro de 2026}

\begin{document}
\maketitle

\begin{abstract}
\noindent
\textbf{Problema.} Para cada documento novo de um caso, decidir se ele é \emph{cópia}, \emph{versão
modificada} ou \emph{documento diferente} em relação a cada documento anterior do mesmo caso.

\smallskip\noindent
\textbf{Representação.} Cada documento vira, uma única vez na ingestão, um texto normalizado
(OCR e layout limpos), dois hashes SHA-256 (texto normalizado e ``esqueleto'' com entidades
mascaradas), uma assinatura MinHash de shingles de 5 tokens e uma sequência de tokens com offsets
no texto. Embeddings foram testados e descartados.

\smallskip\noindent
\textbf{Quase idêntico $\times$ vocabulário parecido.} A comparação usa o \emph{alinhamento} de
tokens, isto é, ordem e continuidade, e não vocabulário. Cada região divergente é classificada como
ruído de OCR, dado variável (nome, valor, data), trecho movido ou conteúdo. Só conteúdo novo faz
de um documento uma versão. Peças da mesma tese escritas por outro escritório não alinham e são
``diferentes''.

\smallskip\noindent
\textbf{Documentos longos.} Hash, MinHash e alinhamento rodam localmente, sem janela de contexto.
O LLM só é chamado na zona cinzenta entre cópia e versão. Ele recebe apenas os trechos
divergentes, no máximo 25, com 8 tokens de contexto cada, cerca de 1.300 tokens por chamada,
qualquer que seja o tamanho do documento.

\smallskip\noindent
\textbf{Avaliação.} Dataset sintético com rótulo exato por construção: 60 casos, 373 documentos e
1.065 pares. Calibração no dev e medição no test. Os documentos são ingeridos na ordem de chegada.
Métricas: F1 por classe, precisão de cópia, taxa de revisão humana, latência e custo.

\smallskip\noindent
\textbf{Resultado.} A cascata completa atinge macro-F1 de 99,4\% e precisão de cópia de 100\%, a
US\$0,0001 por documento, com 5,4\% dos documentos enviados à revisão humana. Embeddings sozinhos
ficam em 65,5\%.

\smallskip\noindent
\textbf{Ação sobre cópia ou versão.} Nada é escondido nem apagado. O documento entra na
\emph{família} do documento casado, como duplicata ou nova versão, com o diff anexado e com os
dados que mudaram em destaque. A zona cinzenta vai para confirmação humana.

\smallskip\noindent
\textbf{Trade-off.} O juiz LLM ganha 0,9 ponto de macro-F1 e elimina falsos ``cópia'', ao custo de
uma p95 de 6,6\,s nos documentos da zona cinzenta. Ele pode rodar de forma assíncrona. Sem ele, o
sistema decide 100\% dos documentos em p95 de 25\,ms, a custo zero de API.
\end{abstract}

\tableofcontents
\bigskip

% ------------------------------------------------------------------------------------------
\section{Contexto e requisitos}

No contencioso de massa, os mesmos escritórios protocolam milhares de petições quase idênticas a
partir de templates, mudando nomes, valores e datas. Três requisitos guiaram o desenho:

\begin{itemize}
  \item \textbf{Os erros têm custos diferentes.} Uma cópia não detectada desperdiça trabalho. Um
    falso ``cópia'' pode esconder um documento que importa. A precisão de cópia é a métrica
    principal.
  \item \textbf{Vocabulário em comum não é sinal.} Peças jurídicas compartilham boilerplate.
  \item \textbf{Volume, custo e auditoria.} Cada decisão precisa ser barata, rápida e rastreável
    até trechos da fonte.
\end{itemize}

O escopo da comparação é o caso: cada documento novo é comparado só com os anteriores do mesmo
caso.

% ------------------------------------------------------------------------------------------
\section{Dataset}

Não há dados de cliente. O dataset é sintético e o rótulo de cada par é exato por construção.

\subsection{Geração}

\begin{enumerate}
  \item \textbf{Templates, gerados uma vez por LLM} (\cfg{gpt-4.1}): boilerplate por escritório,
    conteúdo específico por tese e escritório, uma paráfrase de cada parágrafo e parágrafos
    extras, que são a matéria-prima das versões. Saídas com dado concreto, como R\$, CPF ou data,
    são rejeitadas. Os templates ficam versionados no repositório.
  \item \textbf{Montagem determinística} (seed 13). Cada documento é uma sequência de
    \emph{blocos} com id de origem, mais valores fictícios de entidades e um estilo de
    formatação. Daí saem as mutações: troca de entidades, ruído de PDF e de OCR em três níveis,
    versões leve, média e forte, contestação da ré e negativos difíceis.
  \item \textbf{Rótulo pela proveniência, nunca por métrica de texto.} \emph{Cópia} tem a mesma
    sequência de blocos. \emph{Versão} vem da mesma linhagem de template com outra sequência de
    blocos. \emph{Diferente} vem de linhagens distintas, como outro escritório, outra tese ou
    outro tipo de peça. Assim o gabarito não favorece nenhum método.
\end{enumerate}

\subsection{Composição}

\begin{table}[H]
\centering
\begin{tabular}{lrrr}
\toprule
 & dev & test & total \\
\midrule
casos & 30 & 30 & 60 \\
pares cópia & 92 & 92 & 184 \\
pares versão & 180 & 142 & 322 \\
pares diferente & 256 & 303 & 559 \\
\midrule
pares (total) & 528 & 537 & 1.065 \\
\bottomrule
\end{tabular}
\caption{Pares por rótulo e split. O split é por caso. O test tem 186 documentos. São de 2 a 10
documentos por caso, com mediana de 6. Cada documento tem mediana de 11.521 caracteres e máximo
de 23.624.}
\label{tab:dataset}
\end{table}

Os negativos difíceis existem porque vocabulário parecido não é sinal de cópia:
\begin{itemize}
  \item \cfg{tese_vizinha_mesmo_escritorio}: mesma ré e mesmo boilerplate, com tese vizinha, como
    atraso e cancelamento de voo;
  \item \cfg{mesma_tese_outro_escritorio}: mesma tese e mesmos fatos, com redação de outro
    escritório;
  \item \cfg{outra_tese_mesmo_escritorio}: mesmo boilerplate de consumidor, outro setor;
  \item \cfg{tipo_diferente}: contestação da ré, com as mesmas partes, valores e tema.
\end{itemize}

As superfícies simulam como os documentos chegam à plataforma: 124 documentos digitais, 150 PDFs
exportados, com quebras de linha, paginação e hifenização, e 99 digitalizados com OCR leve, médio
ou forte, com troca de caracteres e palavras coladas ou partidas.

% ------------------------------------------------------------------------------------------
\section{Solução}

\subsection{Representação do documento}

Tudo é calculado uma vez na ingestão e guardado como assinatura (Tabela~\ref{tab:repr}).

\begin{table}[H]
\centering
\small
\begin{tabular}{p{5.6cm}p{8.6cm}}
\toprule
Representação & Uso \\
\midrule
Texto normalizado: OCR/PDF limpo, marcadores de página e cabeçalhos corridos removidos,
parágrafos partidos costurados & Texto canônico. É o que o revisor vê, e todo offset aponta
para ele. \\
SHA-256 do texto normalizado & Cópia exata em $O(1)$. \\
SHA-256 do esqueleto: entidades mascaradas (nomes, CPF/CNPJ, valores, datas, nº CNJ),
minúsculas, sem acento & Cópia com troca de partes, valores e datas, em $O(1)$. \\
Shingles de 5 tokens e MinHash com 128 permutações & Pré-filtro barato. Também serve para LSH
se a busca sair do caso. \\
Tokens mascarados com offsets de caractere & Alinhamento e diff rastreável até a fonte. \\
\bottomrule
\end{tabular}
\caption{Representações por documento.}
\label{tab:repr}
\end{table}

\paragraph{Por que tokens e não parágrafos.} Quebras de página partem parágrafos, e o OCR perde
as linhas em branco e funde o documento inteiro. O parágrafo deixa de ser uma unidade estável. O
token, com offset no texto normalizado, continua estável.

\subsection{Cascata de decisão}

Para cada par formado pelo documento novo e um documento anterior do caso, as etapas vão da mais
barata à mais cara (Figura~\ref{fig:cascata}).

\begin{figure}[H]
\centering
\begin{tikzpicture}[
  node distance=6mm and 10mm,
  box/.style={draw, rounded corners, align=center, minimum width=4.6cm, minimum height=9mm,
              font=\small},
  res/.style={draw, dashed, rounded corners, align=center, font=\small, minimum width=2.4cm},
  >={Stealth}]
  \node[box] (h) {1. Hash normalizado / esqueleto};
  \node[box, below=of h] (m) {2. Pré-filtro MinHash};
  \node[box, below=of m] (a) {3. Alinhamento de tokens\\+ regra calibrada};
  \node[box, below=of a] (j) {4. Juiz LLM\\(só zona cinzenta)};
  \node[res, right=of h] (oc) {cópia};
  \node[res, right=of m] (od) {diferente};
  \node[res, right=of a] (or) {cópia / versão /\\diferente};
  \node[res, right=of j] (oj) {cópia / versão};
  \draw[->] (h) -- node[left, font=\scriptsize]{diferem} (m);
  \draw[->] (m) -- node[left, font=\scriptsize]{acima do limiar} (a);
  \draw[->] (a) -- node[left, font=\scriptsize]{zona cinzenta} (j);
  \draw[->] (h) -- node[above, font=\scriptsize]{igual} (oc);
  \draw[->] (m) -- node[above, font=\scriptsize]{abaixo} (od);
  \draw[->] (a) -- (or);
  \draw[->] (j) -- (oj);
\end{tikzpicture}
\caption{Cascata por par. No test, as etapas decidiram 16 pares por hash, 229 pelo pré-filtro,
272 pela regra e 20 pelo LLM.}
\label{fig:cascata}
\end{figure}

\paragraph{Etapa 3: alinhamento e classificação de regiões.} Os tokens dos dois documentos são
alinhados por Indel/LCS (\texttt{rapidfuzz}). Diferenças separadas por menos de 2 tokens iguais
formam uma \emph{região}, e cada região recebe um tipo:
\begin{itemize}
  \item \textbf{ruído:} similaridade de caracteres acima de 75 com o outro lado, como erro de OCR
    ou hifenização;
  \item \textbf{dado:} pelo menos metade dos tokens são máscaras de entidade, números ou nomes
    próprios;
  \item \textbf{movido:} pelo menos 80\% dos trigramas existem em outro ponto do outro documento;
  \item \textbf{conteúdo:} todo o resto, ou seja, fato, argumento, pedido ou jurisprudência novos.
\end{itemize}
Duas features decidem o rótulo. A \emph{cobertura} é a fração de tokens alinhados. A \emph{razão
de conteúdo} é o número de tokens em regiões de conteúdo dividido pelo tamanho do maior documento.
A regra tem dois limiares, calibrados no dev:
\[
\text{rótulo} =
\begin{cases}
\text{diferente} & \text{se cobertura} < t_{\text{diff}} = 0{,}633,\\
\text{cópia} & \text{se razão de conteúdo} \le t_{\text{cópia}} = 0{,}0113,\\
\text{versão} & \text{caso contrário.}
\end{cases}
\]
$t_{\text{diff}}$ maximiza a acurácia entre ``diferente'' e o resto. $t_{\text{cópia}}$ maximiza
o macro-F1 sob a restrição de precisão de cópia $\ge 98\%$ no dev, porque o falso ``cópia'' é o
erro mais caro. O limiar do pré-filtro MinHash é 0,176: uma fração da menor similaridade
observada entre pares da mesma família no dev.

\paragraph{Etapa 4: juiz LLM na zona cinzenta.} A zona cinzenta é a faixa
\[ t_{\text{cópia}}/2 \;\le\; \text{razão de conteúdo} \;\le\; 2\,t_{\text{cópia}}. \]
Nela, o
\cfg{gpt-4.1-mini}, com o prompt \cfg{q1-judge-v1} e temperatura 0, recebe só as regiões que não
são ruído óbvio, as 25 maiores, cada uma com 8 tokens de contexto. Ele classifica \emph{cada
trecho por id} como ruído de OCR, dado variável, formatação ou conteúdo. A decisão é derivada em
código: versão se e somente se algum trecho for conteúdo. Assim todo veredito aponta para trechos
da fonte, e o campo de decisão livre do modelo é ignorado.

\paragraph{Ação por documento.} O documento novo recebe o rótulo do par mais forte, na ordem
cópia, versão, diferente. Ele é marcado para revisão humana quando esse par caiu na zona cinzenta.

% ------------------------------------------------------------------------------------------
\section{Protocolo experimental}

\begin{itemize}
  \item \textbf{Calibração e teste separados.} Os limiares e a regressão logística são
    ajustados só no dev. Os números são do test: 30 casos, 186 documentos e 537 pares.
  \item \textbf{Simulação da produção.} Os documentos de cada caso de teste são ingeridos na
    ordem de chegada, e cada um é comparado com os anteriores.
  \item \textbf{Métricas por par:} precisão, recall e F1 por classe, macro-F1, matriz de
    confusão e acerto por tipo de relação.
  \item \textbf{Métricas por documento:} acurácia da ação (cópia, versão ou novo), precisão da
    ação ``cópia'' (o documento casado está certo), taxa de revisão humana, latência p50 e p95
    e US\$/documento.
  \item \textbf{Intervalo de confiança} de 95\% por bootstrap com 1.000 reamostragens \emph{sobre
    casos}, porque os pares de um mesmo caso não são independentes.
  \item \textbf{Latência} = CPU local medida + latência original das chamadas de API. O cache
    preserva a latência da primeira chamada.
\end{itemize}

\paragraph{Métodos comparados.}
\begin{itemize}
  \item \cfg{jaccard_bruto}: Jaccard de shingles sobre texto só normalizado.
  \item \cfg{minhash}: MinHash sobre texto limpo e mascarado.
  \item \cfg{embeddings}: cosseno médio de embeddings por trecho (\cfg{text-embedding-3-small}).
  \item \cfg{alinhamento}: hash + alinhamento + regra.
  \item \cfg{logreg}: regressão logística sobre todas as features léxicas, de alinhamento e de
    embeddings.
  \item \cfg{cascata}: hash + pré-filtro + alinhamento + regra.
  \item \cfg{cascata_llm}: cascata + juiz LLM na zona cinzenta.
\end{itemize}
Os baselines de similaridade usam a mesma regra de dois limiares, cada um com sua própria
similaridade, também calibrada no dev.

% ------------------------------------------------------------------------------------------
\section{Resultados}

\begin{table}[H]
\centering
\resizebox{\textwidth}{!}{%
\begin{tabular}{lcccccccc}
\toprule
método & macro-F1 (IC95) & P cópia & R cópia & F1 versão & F1 dif. & revisão & lat. p50/p95 (ms) & US\$/doc \\
\midrule
\cfg{jaccard_bruto} & 69,9 (65,6--73,3) & 94,4 & 18,5 & 78,8 & 100,0 & 0,0 & 6 / 12 & 0 \\
\cfg{minhash} & 73,1 (67,7--78,2) & 85,7 & 26,1 & 79,3 & 100,0 & 0,0 & 9 / 18 & 0 \\
\cfg{embeddings} & 65,5 (60,2--69,9) & 100,0 & 17,4 & 72,3 & 94,6 & 0,0 & 50 / 852 & 0,00008 \\
\cfg{alinhamento} & 98,5 (95,6--100,0) & 98,9 & 95,7 & 98,3 & 100,0 & 0,0 & 13 / 30 & 0 \\
\cfg{logreg} & 98,8 (97,1--100,0) & 95,8 & 100,0 & 98,6 & 100,0 & 2,7 & 55 / 858 & 0,00008 \\
\cfg{cascata} & 98,5 (95,6--100,0) & 98,9 & 95,7 & 98,3 & 100,0 & 0,0 & 11 / 25 & 0 \\
\textbf{\cfg{cascata_llm}} & \textbf{99,4} (97,9--100,0) & \textbf{100,0} & 97,8 & 99,3 & 100,0 & 5,4 & 11 / 6596 & 0,00010 \\
\bottomrule
\end{tabular}}
\caption{Split test, valores em \%. P e R por par. ``Revisão'' é a percentagem de documentos
enviados à revisão humana.}
\label{tab:metodos}
\end{table}

\begin{table}[H]
\centering
\small
\resizebox{\textwidth}{!}{%
\begin{tabular}{lccccc}
\toprule
relação & \cfg{minhash} & \cfg{embeddings} & \cfg{alinhamento} & \cfg{logreg} & \cfg{cascata_llm} \\
\midrule
cópia: ruído de PDF & 100,0 & 100,0 & 100,0 & 100,0 & 100,0 \\
cópia: OCR leve & 57,1 & 0,0 & 100,0 & 100,0 & 100,0 \\
cópia: OCR médio & 0,0 & 0,0 & 100,0 & 100,0 & 100,0 \\
cópia: OCR forte & 0,0 & 0,0 & 91,3 & 100,0 & 91,3 \\
cópia: troca de entidades & 0,0 & 0,0 & 77,8 & 100,0 & 100,0 \\
versão leve & 75,0 & 93,8 & 93,8 & 75,0 & 100,0 \\
versão média & 100,0 & 85,9 & 100,0 & 100,0 & 100,0 \\
versão forte & 100,0 & 80,6 & 100,0 & 100,0 & 100,0 \\
dif.: mesma tese, outro escritório & 100,0 & 93,0 & 100,0 & 100,0 & 100,0 \\
dif.: tese vizinha, mesmo escritório & 100,0 & 100,0 & 100,0 & 100,0 & 100,0 \\
\bottomrule
\end{tabular}}
\caption{Percentagem de pares certos por tipo de relação, no test.}
\label{tab:relacao}
\end{table}

\begin{table}[H]
\centering
\begin{tabular}{lccc}
\toprule
verdade $\backslash$ predição & cópia & versão & diferente \\
\midrule
cópia & 90 & 2 & 0 \\
versão & 0 & 142 & 0 \\
diferente & 0 & 0 & 303 \\
\bottomrule
\end{tabular}
\caption{Matriz de confusão de \cfg{cascata_llm}, por par no test. O juiz recebeu 20 pares, com
26.103 tokens de entrada no total.}
\label{tab:confusao}
\end{table}

\subsection{Ablações}

\begin{table}[H]
\centering
\small
\begin{tabular}{lccp{5.2cm}}
\toprule
ablação & macro-F1 & P / R cópia & leitura \\
\midrule
\cfg{alinhamento} (referência) & 98,5 & 98,9 / 95,7 & --- \\
\cfg{sem_mascara} & 97,0 & 98,8 / 90,2 & troca de dados parece mudança \\
\cfg{sem_layout} & 97,6 & 93,8 / 97,8 & cabeçalho e paginação viram ``conteúdo'' \\
\cfg{alinhamento_sem_hash} & 98,5 & 98,9 / 95,7 & o hash é atalho de custo, não de qualidade \\
\cfg{minhash_boilerplate} & 51,0 & 100,0 / 17,4 & remover o texto repetido entre casos destrói o sinal \\
\bottomrule
\end{tabular}
\caption{Ablações no test, valores em \%. \cfg{minhash_boilerplate} deve ser comparado com
\cfg{minhash}, que tem 73,1.}
\label{tab:ablacoes}
\end{table}

% ------------------------------------------------------------------------------------------
\section{Discussão: respostas às perguntas do case}

\subsection{Como representar cada documento}
Combinação de hashes, MinHash e tokens com offsets (Tabela~\ref{tab:repr}). Texto bruto no LLM foi
descartado por custo, latência, limite de contexto e falta de ponteiro para a fonte.
\textbf{Embeddings foram descartados} pelos números. Eles medem assunto, e o assunto de cópia,
versão e petição da mesma tese de outro escritório é o mesmo. O limiar calibrado para
``diferente'' ficou em cosseno 0,977, e o recall de cópia é 17,4\%: nenhuma cópia com OCR médio,
OCR forte ou troca de entidades foi reconhecida (Tabela~\ref{tab:relacao}). Somados às features
léxicas na \cfg{logreg}, eles não melhoram a precisão de cópia, que fica em 95,8 contra 98,9 da
regra.

\subsection{Quase idêntico $\times$ vocabulário jurídico parecido}
O alinhamento mede ordem e continuidade. Duas petições sobre voo atrasado usam as mesmas palavras,
mas não alinham em sequências longas. Uma cópia alinha quase inteira, mesmo com OCR ruim. Depois do
alinhamento, a regra olha \emph{o que} mudou, e não quanto mudou: uma cópia com todos os dados
trocados continua cópia, e uma versão com um único parágrafo novo vira versão. Todos os negativos
difíceis tiveram 100\% de acerto com alinhamento.

Remover o boilerplate parecia intuitivo e se mostrou o pior erro (Tabela~\ref{tab:ablacoes}). No
contencioso de massa o template \emph{é} o sinal de cópia. Os shingles que se repetem entre casos
são justamente os que marcam a cópia.

\subsection{Documentos longos}
Nenhuma etapa precisa do documento inteiro numa janela de contexto. Hash, MinHash e alinhamento
rodam localmente, em milissegundos (p95 de 25\,ms na \cfg{cascata}). O juiz recebe só os trechos
divergentes, cerca de 1.300 tokens por chamada em média no benchmark, independentemente do tamanho
do documento. \textbf{Limite:} o alinhamento cresce com o produto dos tamanhos, e os documentos do
dataset têm no máximo 23,6 mil caracteres. Para peças de centenas de páginas, a mitigação proposta,
não implementada, é ancorar em blocos com hash igual e alinhar só entre âncoras.

\subsection{Como avaliar se funciona}
Offline, com o protocolo da Seção~4: rótulos por proveniência, split por caso, ingestão em ordem,
IC por bootstrap sobre casos e restrição de precisão de cópia na calibração. Reexecutar o benchmark
lendo o cache de LLM e embeddings reproduz exatamente as mesmas métricas de qualidade. Em produção,
a proposta é:
\begin{itemize}
  \item amostra de pares reais rotulada por advogados, estratificada pela etapa que decidiu, e
    recalibração dos limiares nela;
  \item monitoramento da taxa de ``cópia'' desfeita pelo revisor (falso positivo), da taxa de
    ``novo'' agrupado à mão depois (cópia perdida), da fração na zona cinzenta e da distribuição
    por etapa;
  \item alerta quando a fração de regiões de ruído sobe, sinal de mudança de OCR antes de ela
    virar erro.
\end{itemize}

\subsection{O que fazer com uma cópia ou versão provável}
O sistema nunca esconde nem apaga um documento. Ele agrupa em \emph{famílias}, com um documento
canônico e uma cadeia de versões:
\begin{itemize}
  \item \textbf{Cópia:} entra como duplicata do canônico. As regiões de dado que mudaram, como
    parte, valor ou data, ficam em destaque, porque o valor de uma cópia pode ser outro. O
    trabalho feito no canônico, como classificação, cálculo e minuta de defesa, pode ser
    reaproveitado depois de conferir esses dados.
  \item \textbf{Versão:} entra como nova versão do documento casado, com o diff de conteúdo e
    offsets nos dois textos. O revisor lê só o que mudou.
  \item \textbf{Diferente:} abre uma família nova.
  \item \textbf{Zona cinzenta:} vai para confirmação humana, o que acontece em 5,4\% dos documentos.
  \item \textbf{Auditoria:} cada par guarda a etapa que decidiu, as features, os limiares, as
    regiões e o veredito do juiz por trecho, com modelo e versão do prompt.
\end{itemize}

\subsection{Trade-offs de custo, latência e precisão}
\begin{itemize}
  \item \textbf{Juiz LLM:} +0,9 ponto de macro-F1, troca de entidades de 77,8\% para 100\% e
    precisão de cópia de 100\%. Custa US\$0,00010 por documento, cerca de US\$100 por milhão de
    documentos. A p95 vai de 25\,ms para 6,6\,s, e 5,4\% dos documentos vão para revisão. Como ele
    só atua na zona cinzenta, pode rodar de forma assíncrona: o documento entra com o rótulo da
    regra e é confirmado depois.
  \item \textbf{Embeddings:} mais caros, mais lentos (p95 de 852\,ms, dominada pela API) e
    piores. Ficaram fora.
  \item \textbf{Regressão logística:} recall de cópia de 100\%, mas precisão de cópia de 95,8\%,
    pior no erro mais caro, e menos auditável que dois limiares. A regra ficou.
  \item \textbf{Atalhos:} o hash e o pré-filtro não mudam a qualidade, mas o pré-filtro decide
    229 dos 537 pares sem alinhar.
\end{itemize}

% ------------------------------------------------------------------------------------------
\section{Operação em produção}

\begin{itemize}
  \item \textbf{Escopo por caso.} No dataset, cada caso tem mediana de 6 documentos. O custo por
    documento é dominado pelo alinhamento, em milissegundos.
  \item \textbf{Fila por caso.} A ingestão é serializada dentro do caso, para que documentos que
    chegam juntos se vejam em ordem, e paralela entre casos.
  \item \textbf{Além do caso.} A mesma assinatura MinHash serve para LSH no corpus inteiro. Com
    limiar 0,5 sobre os 373 documentos, o recall de pares da mesma família foi de 92,1\%,
    examinando 21,5 candidatos por documento (5,8\% do corpus), com consulta de 0,02\,ms. Isso
    serve para achar o mesmo modelo de petição em vários processos.
  \item \textbf{Dados sensíveis.} As etapas 1 a 3 rodam dentro da infraestrutura. Só a zona
    cinzenta chama o LLM, e só com trechos divergentes, sob Zero Data Retention.
\end{itemize}

% ------------------------------------------------------------------------------------------
\section{Limitações}

\begin{itemize}
  \item \textbf{Dados sintéticos.} O ruído vem de mutadores determinísticos, e OCR real, anexos e
    tabelas digitalizadas não estão cobertos. Os números ficam perto do teto: o benchmark separa
    bem os métodos, mas superestima a qualidade absoluta esperada em produção.
  \item \textbf{Definição de produto.} A fronteira entre cópia e versão é uma escolha. Hoje
    trocar o valor pedido é ``dado'' e mantém o documento como cópia. Tratar isso como versão
    exige mudar a regra por tipo de região e recalibrar.
  \item \textbf{Documentos muito longos.} O desempenho do alinhamento em peças de centenas de
    páginas não foi medido.
\end{itemize}

% ------------------------------------------------------------------------------------------
\section{Reprodutibilidade}

Os números deste documento vêm de \texttt{results/q1-bench-a0e48b8c85/report.md}. O dataset é a
versão \cfg{q1-synth-v1}, com fingerprint \cfg{d4c0f411b0ec}. O \texttt{run\_id} é o hash de
modelo, versão do prompt, schema, dataset, parâmetros e versão do código.

\begin{verbatim}
uv sync
uv run q1 build                    # dataset (seed 13), sem LLM
uv run q1 bench                    # calibra no dev, mede no test
uv run q1 demo --case-id caso-007  # ação, família e diff de um caso
\end{verbatim}

Toda chamada de LLM e de embeddings passa por um cache endereçado pelo hash de modelo, mensagens,
schema e parâmetros. Reexecutar o benchmark reproduz as mesmas métricas de qualidade e de custo.
Só a latência de CPU local varia, em 1 a 2\,ms.

\end{document}
